\chapter{设计约束与待解决问题}

\section{技术限制}
LingoBridge 平台在目前的研发与试点推广阶段，受到跨国海缆物理通道抖动以及海外静态应用边缘托管代理机制的工程约束，在跨海文件上传与解析维度存在客观的技术限制。

由于系统前端将 PDF 分页解析任务完全下推至学生和教师端本地浏览器执行，在大部分教学课件场景下能实现同步 Canvas 渲染，但当教师尝试上传体积超过 20MB 以上的复杂课件或高容量 Excel 作业包时，边缘反向代理网关易触发 502 或 504 网关超时故障。在当前的系统版本中，平台在前端对上传文件字节数执行了静态拦截；而在系统长期的架构演进路线中，规划采用对象存储直传技术，直接签发带有安全密钥的预签名直传凭证（Presigned URL），使文件数据流直接对接高带宽的云存储服务，从而解决由于代理服务器缓冲区溢出导致的上传超时技术限制。

\section{设计限制与技术风险}
系统目前面临的设计限制与技术风险，主要表现在外部会议平台和自动语音评分接口的企业级准入门槛限制上。

虽然平台在设计上规划了 Microsoft Teams 和腾讯会议 API 的集成，用以自动生成直播房间并同步拉取学生的考勤数据，但由于此类会议系统的企业级开发者账号存在繁琐的企业资质审核流程、漫长的安全性评估周期以及严格的准入门槛限制，导致在试点初期，平台面临无法及时获取腾讯会议及 Microsoft Teams 官方 OAuth 2.0 签名鉴权密钥的技术风险。该风险属于外部非技术控制因素，需要架构团队在系统设计层面提前制定降级预案。

\section{战略 Fallback 降级预案与后续演进}
为了确保在腾讯会议企业开发者官方授权密钥审批延误的情况下，留学生交互课堂的教学闭环仍能跑通，平台架构设计了第三方会议外链分享降级预案。

当系统检测到后端第三方 API 调用凭证缺失或鉴权失败时，系统并不阻塞创建课时节点的核心链路，而是自动切换至降级模式。在此模式下，教师在排课或交互课堂中创建课程时，系统会提供手动的入会链接（Join URL）输入框。教师可以直接将从腾讯会议或 Microsoft Teams 个人客户端中生成的会议邀请入会外链复制并填入其中。

当教师提交该降级数据时，后端网关会调用正则表达式安全解析过滤器，对填入的字符串执行合法性与安全性校验，提取出安全的会议标识。系统前端在加载该课时时，会将该降级链接重新组装，并以高度一致的“进入教室”按钮渲染在学生的交互看板中。学生点击后，可唤醒本地的第三方会议客户端进行授课。此战略降级方案以零 API 授权资费成功达成了核心业务链的跑通，展现了本平台概要设计在应对复杂商业资质瓶颈时的架构柔性。
